\chapter{O Som do Silêncio}

\lettrine{À}{\space medida que nos acalmamos} podemos sentir o som do silêncio na mente. Nós
ouvimo-lo como uma espécie de som de alta frequência, um som a ressoar sempre
presente. Normalmente nunca é percebido. Ora, quando começamos a ouvir aquele
som do silêncio, é um sinal de vazio -- de silêncio da mente. É algo a que
podemos sempre recorrer. À medida que nos concentramos nele e nos dirigimos para
ele, pode tornar-nos bastante pacíficos, com plena felicidade. Meditando nisso,
temos uma forma de deixar que as condições da mente cessem, sem termos que as
suprimir com outra condição. Senão, acabamos a colocar uma condição sobre a
outra.

Este processo é o que normalmente se quer significar por ``kamma''. Por exemplo,
se nos estamos a sentir irritados, começamos a pensar noutra coisa para nos
afastarmos da raiva. Não gostamos do que está a acontecer aqui, então olhamos
para acolá e saímos a correr. Mas se temos uma forma de nos desviarmos dos
fenómenos condicionados para os não-condicionados, então nenhum tipo de kamma
está a ser criado e os hábitos condicionados podem desaparecer e cessar. É como
uma ``válvula de segurança'' na mente, o meio de saída, de forma a que as nossas
% REVISÃO: "formações cármicas" — incoerência terminológica; noutros pontos usa-se "formações kammicas" (de "kamma"). Uniformizar. (Também "\emph{saṅkhāras}" com -s: noutros capítulos é "saṅkhārā".)
formações cármicas (\emph{saṅkhāras}) tenham uma saída, um meio de se esvaziarem
em vez de se recriarem.

Um problema com a meditação é que muitas pessoas a acham monótona. As pessoas
ficam entediadas com o vazio. Querem preencher o vazio com alguma coisa.
Portanto, reconheçamos que mesmo quando a mente está bastante vazia, os desejos
e hábitos ainda estão lá, e eles virão e quererão fazer algo de interessante.
Nós temos que ser pacientes, dispostos a afastarmo-nos do tédio e do desejo de
fazer algo interessante, e contentarmo-nos com o vazio do som do silêncio. E
temos que ser bastante determinados para nos dirigirmos ao som do silêncio.

No entanto, quando começamos a ouvir e compreender melhor a mente, é uma
possibilidade muito realizável para todos nós. Depois de muitos anos de prática,
as formações kammicas grosseiras desaparecem, enquanto as mais subtis começam
também a desaparecer. A mente fica cada vez mais vazia e clara. Mas é preciso
muita paciência, resistência e boa vontade para continuar a praticar em todas as
condições e abandonar até mesmo os nossos pequenos hábitos mais estimados.

Podemos acreditar que o som do silêncio é algo, ou que é uma conquista. No
entanto, não se trata de ter alcançado nada, mas de sabiamente reflectir sobre o
que experienciamos. A forma como se deve reflectir é -- tudo o que vem, acaba
por partir; e a prática consiste em conhecer as coisas como elas são.

Não estou a dar nenhum tipo de identidade -- não há nada para nos apegarmos.
Algumas pessoas, quando ouvem aquele som, querem saber ``É isso a entrada na
corrente?'' Ou ``Temos uma alma?'' Estamos tão agarrados aos conceitos. Tudo o
que podemos saber é que queremos saber algo, queremos ter um rótulo para o nosso
``eu''. Se houver uma dúvida sobre alguma coisa, a dúvida surge e então existe
desejo por alguma coisa. Mas a prática trata-se de desapegar. Continuamos com
aquilo que é, reconhecendo as condições como condições e o incondicionado como
incondicionado. É tão simples como isto.

Até o aspirar religioso é visto como uma condição! Não significa que não devemos
aspirar a, mas sim que simplesmente se deve reconhecer a aspiração em si mesma
como limitada. E o vazio também é não-eu -- apego à ideia de vazio também é
apego. Abandonemos isso! A prática então torna-se a forma de nos afastarmos dos
fenómenos condicionados, não criando nada mais em torno das condições
existentes. Portanto, seja o que for que surja na nossa consciência -- raiva ou
ganância ou qualquer coisa -- reconhecemos que está lá, mas não fazemos nada com
isso. Podemos voltar-nos para o vazio da mente -- para o som do silêncio. Isto
% REVISÃO: pontuação — "oferece às condições como a raiva, uma saída para a cessação, deixamos a condição ir-se embora": o inciso deveria ler-se "oferece às condições, como a raiva, uma saída para a cessação" e a vírgula antes de "deixamos" forma uma frase corrida (usar ponto e vírgula ou "e deixamos").
oferece às condições como a raiva, uma saída para a cessação, deixamos a
condição ir-se embora.

Temos memórias do que fizemos no passado, não temos? Elas surgem na consciência
quando as condições existem para que elas surjam. Isso é o kamma resultante de
se ter feito algo no passado, de se ter agido por ignorância e ter feito coisas
por ganância, ódio, ilusão e assim por diante\ldots{} Quando esse kamma
amadurece no presente, a pessoa ainda tem os impulsos da ganância, ódio e ilusão
que surgem na mente, o kamma resultante. Sempre que agimos sob estes aspectos de
forma ignorante, quando não conscientes, criamos então mais kamma.

As duas vias pelas quais podemos criar kamma são, ou segui-lo ou livrarmo-nos
dele. Quando paramos de fazer isso, os ciclos de kamma têm a oportunidade de
cessar. O kamma resultante que surgir tem um escape, uma ``saída de emergência''
para a cessação.
